Through comprehensive analysis of 43 states across 645 optimization configurations, we establish a statistically validated \textbf{42\% effectiveness threshold} for near-proportional VRA compliance through algorithmic redistricting. States above this threshold achieve 80-114\% of proportional MM targets (62\% average success rate, 92\% achieve targets), while states below 37\% achieve 0-67\% of targets (0\% success rate, 0\% achieve targets). The threshold pattern holds consistently across diverse geographic regions, demographic compositions, and state sizes (r=0.78, p$<$1e-08).

\subsection{Core Contributions}

\textbf{Comprehensive national validation}: First systematic empirical determination of VRA effectiveness threshold using comprehensive sample (N=43 states, 645 configurations) with statistical validation (r=0.78, t=7.39, p$<$1e-08). The 42\% threshold operationalizes Thornburg v. Gingles' ``sufficiently large'' requirement.

\textbf{Threshold taxonomy}: Distinguish three concepts: feasibility ($\sim$2\%, can create any MM districts), effectiveness (42\%, can achieve near-proportional representation), and proportionality ($\sim$80\%, reliably achieve full proportional representation). VRA Section 2 requires effectiveness, not merely feasibility.

\textbf{Geographic generalizability}: Demonstrate threshold consistency across all US regions, demographic compositions (Black, Hispanic, Asian, multi-racial), and state sizes (2-52 districts). Pattern reflects fundamental demographic-geographic constraints, not regional artifacts.

\textbf{Evidence-based policy framework}: Provide courts, legislatures, and redistricting commissions with empirically validated guidance for VRA compliance assessment, reducing litigation uncertainty and enabling ex ante feasibility evaluation.

\subsection{Practical Implications}

\textbf{Courts}: Can objectively evaluate Gingles prong 1 using 42\% threshold. States above threshold have presumptive feasibility; states below require extraordinary geographic clustering evidence.

\textbf{Legislatures}: Understand VRA obligations before map drawing. States above threshold should proactively create MM districts; states below can document infeasibility.

\textbf{Advocates}: Focus resources on states above 42\% where compliance clearly achievable. In borderline states (37--42\%), commission detailed geographic clustering analysis.

\subsection{Theoretical Significance}

The 42\% threshold emerges from fundamental mathematical constraints, not algorithmic artifacts:

\begin{itemize}
    \item \textbf{Proportionality}: Cannot create $X$\% MM districts from $<X$\% minority population
    \item \textbf{Geographic scatter}: Distributing minority population across $k$ districts while maintaining contiguity requires sufficient state-wide base
    \item \textbf{50\% threshold}: Creating majority-minority districts (not plurality-minority) requires concentrating enough minority population in each target district
\end{itemize}

Future algorithmic improvements may reduce threshold modestly but cannot eliminate it entirely without violating geographic constraints (contiguity, compactness).

\subsection{Closing Remarks}

The Voting Rights Act mandates creation of majority-minority districts ``where feasible,'' yet ``feasible'' has remained operationally undefined for decades, resolved only through case-by-case litigation. Our 43-state analysis establishes that effectiveness for near-proportional representation correlates strongly with state-level minority percentage, with 42\% serving as the validated threshold (r=0.78, p$<$1e-08).

This comprehensive empirical framework transforms feasibility assessment from subjective expert testimony to quantitatively-informed analysis. Courts can evaluate Gingles prong 1 claims with empirical evidence as contextual guidance, legislatures can assess where proportional MM representation is achievable ex ante, and redistricting commissions can optimize for MM district creation where geographically feasible. Importantly, courts will continue to evaluate district-by-district feasibility based on case-specific factors; our state-level threshold provides empirical context for these individualized determinations, not a dispositive rule.

The 42\% effectiveness threshold represents not algorithmic limitation but fundamental demographic-geographic reality: below this level, edge-weighted optimization can create MM districts but cannot achieve near-proportional representation while maintaining contiguity and compactness. Distinguishing partial effectiveness from high effectiveness prevents both mandating the impossible and failing to realize genuine opportunities for proportional minority representation where geography permits.

Our findings establish empirical foundation for the next generation of VRA compliance assessment, enabling evidence-based decisions about where proportional representation is achievable and where geographic constraints impose genuine limits.
